\section{第三条主线：Benchmark、Environment 与 RL 共演化}

\subsection{没有 solid benchmark，就很难形成统一范式}
panel 里对 benchmark 的讨论很值得单独拎出来。嘉宾们并不满足于 ``有几个任务集''，而是强调 benchmark 必须足够 solid，能够覆盖不同 agent 能力、环境反馈和评测维度。只有 benchmark 广度与难度足够，研究社区才能判断当前方法到底改进了什么、哪里还失效、哪些能力只是过拟合在小范围任务上。

\begin{importantbox}{Benchmark 在 Agent 领域承担的角色比在纯文本任务里更重}
原因在于 Agent 涉及的变量更多：
\begin{itemize}
  \item 输入不再只是 prompt，还包括环境状态；
  \item 输出不再只是文本，还包括 action sequence；
  \item 评测不再只看答案，还要看过程是否稳定、可复现、可恢复；
  \item 成功与失败常常受 environment fidelity 影响。
\end{itemize}
\end{importantbox}

\subsection{Environment 是训练对象，也是评测基础设施}
多位嘉宾在讨论 RL 时都把 environment 放在非常高的位置。环境不是一个背景板，而是 reward、探索空间、tool space 与 benchmark 难度的共同来源。尤其在 GUI、mobile agent、computer use 这类任务中，如果环境是静态的、伪造的或不可复现的，那么训练和评测都会失真。

\begin{knowledgebox}{为什么 Agent 研究会越来越像 ``环境工程''}
因为 scaling of acting 不够，必须继续 scaling of environment。讨论中明确提到了：
\begin{itemize}
  \item 需要更多 infra、更多任务环境、更多数据；
  \item 需要更大 tool space 与 action space；
  \item 需要真实环境、云手机/云电脑或可复现的沙箱；
  \item 需要 world model 或 environment model 来降低构造成本。
\end{itemize}
\end{knowledgebox}

\subsection{Self-evolving 更像方法论，不是单一任务}
关于 self-evolving，嘉宾指出了一个很容易被忽略的点：它不像 reasoning 那样一眼就对应到某个具体能力，更像一套围绕不同任务自我改进的方法论。因此 code、math、GUI、视觉等领域很难立刻共享完全统一的范式。短期内更现实的路径，是在各个子领域先形成局部共识，再逐渐向更广泛的统一方法推进。

\begin{warningbox}{不要期待所有 Agent 任务立刻共用一套训练 recipe}
文本、代码、GUI、视觉的 environment 差异太大，奖励密度、行动空间和失败模式都不同。强行追求统一配方，很容易把关键细节抹平，得到一个对谁都不够好的折中方案。
\end{warningbox}

\subsection{本章小结}
Agent 研究想走向统一范式，绕不开 benchmark 和 environment。它们既决定 RL 是否可训练，也决定评测是否可信。相比 ``再找一个更好的 prompt''，这条主线更像是在建设一个能持续迭代 Agent 能力的科研基础设施。

